\chapter{Simple Linear Regression and Correlation}

\begin{nuance}{how to use this chapter}
This chapter follows Chapter~11 of the course textbook section by section, with the same titles and order. Regression builds an empirical model for how a response $Y$ changes with a regressor $x$, then uses the machinery of Chapters~2 and~3 (confidence intervals and tests) on the fitted coefficients. The worked example below is carried through every section so that all the formulas can be checked against one set of numbers. Numbering follows the textbook with the chapter number lowered by $6$: Section $5.2$ of this book is Section $11$-$2$ of the textbook, and Section $5.2.1$ is its Section $11$-$2.1$.
\end{nuance}

\section{Empirical Models}

Many problems involve the relationship between two or more variables: the strength of a weld and its temperature, the fuel use of a car and its weight. An \emph{empirical model} is an equation fitted to observed data that describes this relationship without being derived from first principles. Typically the response $Y$ is random and the regressor $x$ is a quantity whose value we choose or measure with negligible error.

\begin{definition}[Simple linear regression model]
The \emph{simple linear regression model} is
\[ Y=\beta_0+\beta_1x+\varepsilon, \]
where $\beta_0$ (the intercept) and $\beta_1$ (the slope) are unknown \emph{regression coefficients} and $\varepsilon$ is a random error with $\Exp\varepsilon=0$ and $\Var\varepsilon=\sigma^2$. For data $(x_i,y_i)$, $i=1,\dots,n$: $Y_i=\beta_0+\beta_1x_i+\varepsilon_i$ with the $\varepsilon_i$ independent. For inference we also assume the errors are normal, $\varepsilon_i\sim\Normal(0,\sigma^2)$.
\end{definition}

Then for fixed $x$, $\Exp(Y\mid x)=\beta_0+\beta_1x$ (the mean response lies on a straight line), $\Var(Y\mid x)=\sigma^2$ (constant variance), and $\beta_1$ is the change in the mean response per unit change in $x$.

\begin{example}\label{ex:regdata}
Throughout this chapter we use the following $n=8$ observations.
\begin{center}
\begin{tabular}{c|cccccccc}
\toprule
$x$ & 1 & 2 & 3 & 4 & 5 & 6 & 7 & 8\\
$y$ & 2.1 & 3.9 & 6.2 & 7.8 & 10.1 & 12.2 & 13.8 & 16.1\\
\bottomrule
\end{tabular}
\end{center}
\end{example}

\section{Simple Linear Regression}

\begin{figure}[h]
\centering
\begin{tikzpicture}[scale=0.55]
  \draw[axis] (0,0)--(9.4,0) node[right,lab]{$x$}; \draw[axis] (0,0)--(0,9.4) node[above,lab]{$y$};
  \draw[thick,accent2] (0.3,0.35)--(9,9.0);
  \foreach \x/\y in {1/1.05,2/1.95,3/3.1,4/3.9,5/5.05,6/6.1,7/6.9,8/8.05}{
    \node[pt] at (\x,\y){}; \draw[guide] (\x,\y)--(\x,{0.0357*0.5+\x*0.99881}); }
\end{tikzpicture}
\caption{The least squares line minimises the sum of squared vertical distances (dashed) from the data points to the line.}
\end{figure}

\subsection*{The method of least squares}

We estimate $\beta_0,\beta_1$ by choosing the line that makes the data as close as possible, in the sense of minimising the sum of squared vertical deviations
\[ L(\beta_0,\beta_1)=\sum_{i=1}^n\big(y_i-\beta_0-\beta_1x_i\big)^2. \]
Setting $\partial L/\partial\beta_0=0$ and $\partial L/\partial\beta_1=0$ gives the \emph{least squares normal equations}
\[ n\hat\beta_0+\hat\beta_1\sum x_i=\sum y_i,\qquad\hat\beta_0\sum x_i+\hat\beta_1\sum x_i^2=\sum x_iy_i. \]

\begin{definition}[Corrected sums of squares]
$S_{xx}=\sum(x_i-\bar{x})^2=\sum x_i^2-\dfrac{(\sum x_i)^2}n$, $\ S_{xy}=\sum(x_i-\bar{x})(y_i-\bar{y})=\sum x_iy_i-\dfrac{(\sum x_i)(\sum y_i)}n$, and $S_{yy}=\sum(y_i-\bar{y})^2$.
\end{definition}

\begin{theorem}[Least squares estimates]\label{thm:ls}
\[ \hat\beta_1=\frac{S_{xy}}{S_{xx}},\qquad\hat\beta_0=\bar{y}-\hat\beta_1\bar{x}. \]
The fitted line is $\hat{y}=\hat\beta_0+\hat\beta_1x$ and the \emph{residuals} are $e_i=y_i-\hat{y}_i$. The line passes through $(\bar{x},\bar{y})$.
\end{theorem}

\begin{proof}
The first normal equation gives $\hat\beta_0=\bar{y}-\hat\beta_1\bar{x}$. Substituting into the second: $\hat\beta_1\big(\sum x_i^2-n\bar{x}^2\big)=\sum x_iy_i-n\bar{x}\bar{y}$, i.e.\ $\hat\beta_1S_{xx}=S_{xy}$. The second derivatives show that this is a minimum ($L$ is a convex quadratic).
\end{proof}

\begin{example}[Continuing Example~\ref{ex:regdata}]\label{ex:regfit}
$\bar{x}=4.5$, $\bar{y}=9.025$, $S_{xx}=42$, $\sum xy=408.8$ so $S_{xy}=408.8-8(4.5)(9.025)=83.9$ and $S_{yy}=819.4-8(9.025)^2=167.795$. Hence $\hat\beta_1=83.9/42=1.9976$ and $\hat\beta_0=9.025-1.9976(4.5)=0.0357$. The fitted line is $\hat{y}=0.0357+1.9976\,x$: each unit increase in $x$ raises the mean response by about $2.0$.
\end{example}

\subsection*{Estimating $\sigma^2$}

The error variance is estimated from the residuals.

\begin{definition}[Error sum of squares, $\hat\sigma^2$]
$SS_E=\sum_{i=1}^ne_i^2=\sum(y_i-\hat{y}_i)^2$, and
\[ \hat\sigma^2=MS_E=\frac{SS_E}{n-2}. \]
A convenient computing formula is $SS_E=S_{yy}-\hat\beta_1S_{xy}$.
\end{definition}

The divisor $n-2$ is the degrees of freedom: two parameters ($\beta_0,\beta_1$) were estimated from the data, so two degrees of freedom are lost, and with it $\Exp[\hat\sigma^2]=\sigma^2$ (shown below). In Example~\ref{ex:regfit}: $SS_E=167.795-1.9976(83.9)=0.195$ and $\hat\sigma^2=0.195/6=0.0325$, so $\hat\sigma=0.180$.

\section{Properties of the Least Squares Estimators}

Write $\hat\beta_1=\sum c_iY_i$ with $c_i=(x_i-\bar{x})/S_{xx}$; the $x_i$ are fixed, so $\hat\beta_1$ is a linear combination of the independent $Y_i$.

\begin{theorem}[Sampling properties]\label{thm:lsprops}
Under the model assumptions:
\begin{enumerate}[label=(\roman*)]
  \item $\Exp\hat\beta_1=\beta_1$ and $\Var\hat\beta_1=\sigma^2/S_{xx}$;
  \item $\Exp\hat\beta_0=\beta_0$ and $\Var\hat\beta_0=\sigma^2\Big(\dfrac1n+\dfrac{\bar{x}^2}{S_{xx}}\Big)$;
  \item $\Exp\hat\sigma^2=\sigma^2$;
  \item if the errors are normal, $\hat\beta_0,\hat\beta_1$ are normal, and $(n-2)\hat\sigma^2/\sigma^2\sim\chi^2_{n-2}$ independently of $(\hat\beta_0,\hat\beta_1)$.
\end{enumerate}
The estimated standard errors are $\se(\hat\beta_1)=\sqrt{\hat\sigma^2/S_{xx}}$ and $\se(\hat\beta_0)=\sqrt{\hat\sigma^2(1/n+\bar{x}^2/S_{xx})}$.
\end{theorem}

\begin{proof}
(i) Since $\sum c_i=0$ and $\sum c_ix_i=1$: $\Exp\hat\beta_1=\sum c_i(\beta_0+\beta_1x_i)=\beta_1$, and $\Var\hat\beta_1=\sigma^2\sum c_i^2=\sigma^2/S_{xx}$. (ii) $\hat\beta_0=\bar{Y}-\hat\beta_1\bar{x}$; the terms $\bar{Y}$ and $\hat\beta_1$ are uncorrelated (because $\sum c_i=0$), so $\Var\hat\beta_0=\sigma^2/n+\bar{x}^2\sigma^2/S_{xx}$. (iii) and (iv) are proved in the matrix setting in Chapter~6 (Theorem~\ref{thm:mlr-dist}): the residual vector lies in an $(n-2)$-dimensional space.
\end{proof}

\begin{theorem}[Gauss--Markov]\label{thm:gm1}
Among all \emph{linear unbiased} estimators of $\beta_1$ (estimators of the form $\sum a_iY_i$), $\hat\beta_1$ has the smallest variance; likewise for $\hat\beta_0$. The least squares estimators are the \emph{best linear unbiased estimators} (BLUE), requiring only uncorrelated errors with equal variance, not normality.
\end{theorem}

\begin{interview}{what makes the slope precise?}
$\Var\hat\beta_1=\sigma^2/S_{xx}$ shrinks when the noise $\sigma$ is small, when $n$ is large, and when the $x$-values are \emph{spread out} (large $S_{xx}$). A designed experiment should place regressor values at the extremes of the range: half at each end maximises $S_{xx}$ for given $n$ and range (though it cannot detect curvature, so a few middle points are added). Extrapolating beyond the range of the data is dangerous, since the fitted line says nothing about whether the relationship stays linear there.
\end{interview}

\section{Some Comments on Uses of Regression (CD Only)}

Regression is used for \emph{description} (summarising a trend), \emph{prediction} (estimating the response at a new $x$ within the range of the data) and \emph{control} (choosing $x$ to achieve a target $y$). Cautions: (1) a fitted line is valid only in the region of the observed $x$, (2) a strong fit shows association, not causation, (3) the regression of $y$ on $x$ is not the regression of $x$ on $y$ (the two lines differ unless the correlation is $\pm1$), (4) the model assumes $x$ is measured without error, and (5) a few influential points can control the whole fit.

\section{Hypothesis Tests in Simple Linear Regression}

\subsection{Use of $t$-Tests}

With normal errors, standardising each coefficient with its estimated standard error gives a $t$ statistic.

\begin{theorem}[$t$-tests on the coefficients]\label{thm:regt}
$T=\dfrac{\hat\beta_1-\beta_{1,0}}{\se(\hat\beta_1)}\sim t_{n-2}$ when $\beta_1=\beta_{1,0}$, and similarly $T=\dfrac{\hat\beta_0-\beta_{0,0}}{\se(\hat\beta_0)}\sim t_{n-2}$. To test $H_0:\beta_1=\beta_{1,0}$ against $H_1:\beta_1\ne\beta_{1,0}$, reject if $\abs{t_0}>t_{\alpha/2,n-2}$.
\end{theorem}

The most important special case is $H_0:\beta_1=0$ against $H_1:\beta_1\ne0$, the \emph{test for significance of regression}. Failing to reject means there is no evidence of a linear relationship between $x$ and $Y$ (either $x$ is of little value in explaining $Y$, or the true relationship is not linear); rejecting means $x$ is useful in explaining $Y$ (the line is a better model than a constant, though a curve might be better still).

\begin{example}
For Example~\ref{ex:regfit}: $\se(\hat\beta_1)=\sqrt{0.0325/42}=0.0278$ and $t_0=1.9976/0.0278=71.8$ with $6$ degrees of freedom, far beyond $t_{0.005,6}=3.707$: reject $H_0:\beta_1=0$ ($P<0.0001$). For the intercept, $\se(\hat\beta_0)=\sqrt{0.0325(1/8+20.25/42)}=0.1405$ and $t_0=0.0357/0.1405=0.25$: no evidence that the intercept differs from $0$.
\end{example}

\subsection{Analysis of Variance Approach to Test Significance of Regression}

The total variability of the responses splits into a part explained by the regression and a part left over.

\begin{theorem}[Analysis of variance identity]\label{thm:anova-reg}
\[ \underbrace{\sum(y_i-\bar{y})^2}_{SS_T=S_{yy}}=\underbrace{\sum(\hat{y}_i-\bar{y})^2}_{SS_R=\hat\beta_1S_{xy}}+\underbrace{\sum(y_i-\hat{y}_i)^2}_{SS_E}, \]
with degrees of freedom $n-1=1+(n-2)$.
\end{theorem}

\begin{proof}
Write $y_i-\bar{y}=(\hat{y}_i-\bar{y})+(y_i-\hat{y}_i)$ and square; the cross term is $2\sum(\hat{y}_i-\bar{y})e_i=2\hat\beta_1\sum(x_i-\bar{x})e_i=0$, because the normal equations make the residuals sum to zero and orthogonal to the $x_i$.
\end{proof}

The ratio of mean squares gives an $F$ test equivalent to the $t$-test.

\begin{theorem}[$F$-test for significance of regression]\label{thm:regF}
Let $MS_R=SS_R/1$ and $MS_E=SS_E/(n-2)$. Under $H_0:\beta_1=0$, $F_0=MS_R/MS_E\sim F_{1,n-2}$, and $H_0$ is rejected if $f_0>f_{\alpha,1,n-2}$. Moreover $F_0=T_0^2$, where $T_0=\hat\beta_1/\se(\hat\beta_1)$.
\end{theorem}

\begin{center}
\begin{tabular}{lcccc}
\toprule
Source & Sum of squares & df & Mean square & $F_0$ \\
\midrule
Regression & $SS_R$ & $1$ & $MS_R$ & $MS_R/MS_E$ \\
Error & $SS_E$ & $n-2$ & $MS_E=\hat\sigma^2$ & \\
Total & $SS_T$ & $n-1$ & & \\
\bottomrule
\end{tabular}
\end{center}

\begin{example}
$SS_T=167.795$, $SS_R=1.9976(83.9)=167.600$, $SS_E=0.195$; $f_0=167.600/0.0325=5157$, which equals $t_0^2=71.8^2$ and exceeds $f_{0.05,1,6}=5.99$.
\end{example}

\section{Confidence Intervals}

\subsection{Confidence Intervals on the Slope and Intercept}

\begin{theorem}[Intervals for the coefficients]\label{thm:regci}
A $100(1-\alpha)\%$ confidence interval on $\beta_1$ is $\hat\beta_1\pm t_{\alpha/2,n-2}\sqrt{\hat\sigma^2/S_{xx}}$, and on $\beta_0$ is $\hat\beta_0\pm t_{\alpha/2,n-2}\sqrt{\hat\sigma^2(1/n+\bar{x}^2/S_{xx})}$. A confidence interval on $\sigma^2$ is $\dfrac{(n-2)\hat\sigma^2}{\chi^2_{\alpha/2,n-2}}\le\sigma^2\le\dfrac{(n-2)\hat\sigma^2}{\chi^2_{1-\alpha/2,n-2}}$.
\end{theorem}

\begin{example}
For the slope: $1.9976\pm2.447(0.0278)=1.9976\pm0.068$, i.e.\ $1.930\le\beta_1\le2.066$.
\end{example}

\subsection{Confidence Interval on the Mean Response}

The \emph{mean response} at $x=x_0$ is $\mu_{Y|x_0}=\beta_0+\beta_1x_0$, estimated by $\hat\mu_{Y|x_0}=\hat\beta_0+\hat\beta_1x_0$, which is normal and unbiased with variance $\sigma^2\Big[\dfrac1n+\dfrac{(x_0-\bar{x})^2}{S_{xx}}\Big]$.

\begin{theorem}[Interval on the mean response]\label{thm:meanresp}
A $100(1-\alpha)\%$ confidence interval on $\mu_{Y|x_0}$ is
\[ \hat\mu_{Y|x_0}\pm t_{\alpha/2,n-2}\sqrt{\hat\sigma^2\Big[\frac1n+\frac{(x_0-\bar{x})^2}{S_{xx}}\Big]}. \]
\end{theorem}

The interval is narrowest at $x_0=\bar{x}$ and widens as $x_0$ moves away from the centre of the data: our knowledge of the line is best in the middle.

\begin{example}
At $x_0=5$: $\hat\mu=0.0357+1.9976(5)=10.024$; the standard error is $\sqrt{0.0325(1/8+0.25/42)}=0.0652$, so the $95\%$ interval is $10.024\pm2.447(0.0652)=10.024\pm0.160$.
\end{example}

\section{Prediction of New Observations}

A new individual observation $Y_0$ at $x_0$ equals the mean response plus a new error $\varepsilon_0$ independent of the data.

\begin{theorem}[Prediction interval]\label{thm:regpi}
A $100(1-\alpha)\%$ prediction interval for a future observation $Y_0$ at $x_0$ is
\[ \hat{y}_0\pm t_{\alpha/2,n-2}\sqrt{\hat\sigma^2\Big[1+\frac1n+\frac{(x_0-\bar{x})^2}{S_{xx}}\Big]}. \]
\end{theorem}

\begin{proof}
$Y_0-\hat{Y}_0$ has mean $0$ and variance $\sigma^2+\sigma^2[1/n+(x_0-\bar{x})^2/S_{xx}]$ (the two parts are independent), and standardising with $\hat\sigma$ gives a $t_{n-2}$ pivot.
\end{proof}

\begin{example}
At $x_0=5$: $\sqrt{0.0325(1+0.1310)}=0.1917$, so the prediction interval is $10.024\pm2.447(0.1917)=10.024\pm0.469$, much wider than the interval $\pm0.160$ for the mean response: even with the line known exactly, a single new observation still scatters by $\sigma$.
\end{example}

\section{Adequacy of the Regression Model}

The tests and intervals are valid only if the model assumptions hold: linearity, independent errors with constant variance and (for small $n$) normality.

\subsection{Residual Analysis}

\begin{definition}[Residuals]
$e_i=y_i-\hat{y}_i$. The \emph{standardised residuals} are $d_i=e_i/\hat\sigma$; about $95\%$ of them should lie within $\pm2$ if the model is correct.
\end{definition}

Graphical checks:
\begin{itemize}
  \item \textbf{Normal probability plot of residuals:} points near a straight line support normality; curvature indicates skew or heavy tails.
  \item \textbf{Residuals versus fitted values (or versus $x$):} a random horizontal band is good; a curved pattern means the straight line is the wrong shape (try a quadratic term or a transformation); a funnel (spread increasing with $\hat{y}$) indicates non-constant variance (try transforming $y$ or weighted least squares).
  \item \textbf{Residuals versus time order:} runs of same-sign residuals indicate correlated errors.
\end{itemize}

\subsection{Coefficient of Determination ($R^2$)}\label{sec:regr2}

\begin{definition}[$R^2$]
$R^2=\dfrac{SS_R}{SS_T}=1-\dfrac{SS_E}{SS_T}$, the proportion of the variability in $y$ explained by the regression, with $0\le R^2\le1$.
\end{definition}

For the data above $R^2=167.600/167.795=0.9988$. Cautions: a large $R^2$ does not prove that the linear model is correct (a curved relationship over a limited range can give $R^2$ near $1$), and a small $R^2$ is not necessarily a bad model if the noise is naturally large. $R^2$ never decreases when regressors are added, so it cannot be used to compare models with different numbers of regressors (see Chapter~6).

\subsection{Lack-of-Fit Test (CD Only)}

If there are \emph{repeated observations} at some values of $x$ (replicates), the error sum of squares can be split into pure error (variation among replicates) and lack of fit:
\[ SS_E=SS_{PE}+SS_{LOF},\qquad SS_{PE}=\sum_i\sum_j(y_{ij}-\bar{y}_i)^2\ (n-m\text{ df}),\quad SS_{LOF}=SS_E-SS_{PE}\ (m-2\text{ df}), \]
where $m$ is the number of distinct $x$-values. Under the hypothesis that the straight line is adequate, $F_0=\dfrac{SS_{LOF}/(m-2)}{SS_{PE}/(n-m)}\sim F_{m-2,n-m}$; a large $F_0$ is evidence of curvature.

\section{Transformations to a Straight Line}

If a scatter plot is clearly curved, a transformation of $x$, $y$ or both may make the relationship linear, after which the methods above apply.

\begin{center}
\begin{tabular}{llll}
\toprule
Nonlinear model & Transformation & Linear form & Back-transform \\
\midrule
$y=\beta_0e^{\beta_1x}$ & $y^*=\ln y$ & $y^*=\ln\beta_0+\beta_1x$ & $\beta_0=e^{\text{intercept}}$ \\
$y=\beta_0x^{\beta_1}$ & $y^*=\ln y,\ x^*=\ln x$ & $y^*=\ln\beta_0+\beta_1x^*$ & $\beta_0=e^{\text{intercept}}$ \\
$y=\beta_0+\beta_1\ln x$ & $x^*=\ln x$ & $y=\beta_0+\beta_1x^*$ & --- \\
$y=1/(\beta_0+\beta_1x)$ & $y^*=1/y$ & $y^*=\beta_0+\beta_1x$ & --- \\
$y=\dfrac x{\beta_0x-\beta_1}$ & $y^*=1/y,\ x^*=1/x$ & $y^*=\beta_0-\beta_1x^*$ & --- \\
\bottomrule
\end{tabular}
\end{center}

\section{More About Transformations (CD Only)}

Two cautions. First, a transformation changes the error structure along with the mean: if $y=\beta_0e^{\beta_1x}\varepsilon$ with a multiplicative error $\varepsilon$, then $\ln y$ has an additive error and the log-linear fit is correct; if the original error is additive the log-transformed errors are not constant-variance. Always re-check the residuals \emph{after} transforming. Second, predictions made on the transformed scale (a prediction of $\ln y$) must be back-transformed, and the back-transformed mean of $\ln y$ estimates the \emph{median} of $y$, not its mean. For stabilising variance a common choice is $\sqrt y$ for counts and $\ln y$ when the standard deviation is proportional to the mean; the Box--Cox family $y^{(\lambda)}=(y^\lambda-1)/\lambda$ chooses $\lambda$ by maximum likelihood.

\section{Correlation}

So far $x$ was a fixed regressor. Suppose instead that $(X,Y)$ are \emph{both} random, with a bivariate normal distribution, and we are interested in their strength of linear association.

\begin{definition}[Correlation coefficient]
The population correlation is $\rho=\dfrac{\Cov(X,Y)}{\sigma_X\sigma_Y}\in[-1,1]$. For a sample the \emph{sample correlation coefficient} is
\[ R=\frac{S_{xy}}{\sqrt{S_{xx}S_{yy}}}=\hat\beta_1\sqrt{\frac{S_{xx}}{S_{yy}}}. \]
\end{definition}

Relations: $R^2$ (the square of $R$) equals the coefficient of determination of Section~\ref{sec:regr2}; $R$ has the same sign as $\hat\beta_1$; $R=\pm1$ if and only if all points lie exactly on a line; and $\hat\beta_1=R\,s_y/s_x$.

\begin{theorem}[Test and interval for $\rho$]\label{thm:corr}
If $(X,Y)$ is bivariate normal:
\begin{enumerate}[label=(\roman*)]
  \item to test $H_0:\rho=0$ use $T_0=\dfrac{R\sqrt{n-2}}{\sqrt{1-R^2}}\sim t_{n-2}$ under $H_0$ (identical to the $t$-test of $H_0:\beta_1=0$);
  \item for other hypotheses and intervals use Fisher's transformation: $Z=\tfrac12\ln\dfrac{1+R}{1-R}=\operatorname{arctanh}R$ is approximately $\Normal\Big(\operatorname{arctanh}\rho,\ \dfrac1{n-3}\Big)$. A $100(1-\alpha)\%$ interval for $\rho$ is $\tanh\big(\operatorname{arctanh}r\pm z_{\alpha/2}/\sqrt{n-3}\big)$.
\end{enumerate}
\end{theorem}

\begin{example}
For the example data $r=\sqrt{0.9988}=0.9994$; $t_0=0.9994\sqrt6/\sqrt{1-0.9988}=70.7$ (equals the regression $t$ up to rounding). For an interval with $n=20$ and $r=0.6$: $\operatorname{arctanh}0.6=0.693$, so $0.693\pm1.96/\sqrt{17}=0.693\pm0.475$, and $\rho\in(\tanh0.218,\tanh1.168)=(0.21,0.82)$.
\end{example}

\begin{pitfall}{correlation is not causation, and zero correlation is not independence}
A high correlation can come from a common cause or pure coincidence. Correlation measures \emph{linear} association only: $Y=X^2$ with $X$ symmetric about zero has $\rho=0$ yet $Y$ is a deterministic function of $X$. And a single outlier can create or destroy a large correlation; always look at the scatter plot (Anscombe's quartet has four very different data sets with identical $r=0.816$ and identical regression lines).
\end{pitfall}


\begin{summarybox}
\begin{itemize}
  \item Model $Y=\beta_0+\beta_1x+\varepsilon$, $\varepsilon\sim\Normal(0,\sigma^2)$. Least squares: $\hat\beta_1=S_{xy}/S_{xx}$, $\hat\beta_0=\bar{y}-\hat\beta_1\bar{x}$, $\hat\sigma^2=SS_E/(n-2)$ with $SS_E=S_{yy}-\hat\beta_1S_{xy}$.
  \item $\Var\hat\beta_1=\sigma^2/S_{xx}$, $\Var\hat\beta_0=\sigma^2(1/n+\bar{x}^2/S_{xx})$; $\hat\beta$'s are BLUE (Gauss--Markov).
  \item Tests and intervals use $t_{n-2}$; $H_0:\beta_1=0$ is also tested by $F_0=MS_R/MS_E=t_0^2$ with $SS_T=SS_R+SS_E$.
  \item Mean response: $\hat{y}_0\pm t\sqrt{\hat\sigma^2[1/n+(x_0-\bar{x})^2/S_{xx}]}$; new observation: add $1$ inside the bracket; both are narrowest at $\bar{x}$.
  \item Check the residuals (normal plot, versus $\hat{y}$); $R^2=SS_R/SS_T$; lack of fit needs replicates; transformations linearise some nonlinear models.
  \item $r=S_{xy}/\sqrt{S_{xx}S_{yy}}$, $t_0=r\sqrt{n-2}/\sqrt{1-r^2}$; Fisher's $z$ gives intervals for $\rho$.
\end{itemize}
\end{summarybox}

\section{Exercises}

\begin{exercise}
For the data $x$: $2,4,5,7,9,10$ and $y$: $5.1,8.3,9.8,14.2,18.0,19.6$, find the least squares line, $\hat\sigma^2$, the $95\%$ confidence interval for the slope, and test $H_0:\beta_1=2$ against $\beta_1\ne2$.
\end{exercise}

\begin{exercise}
Construct the analysis of variance table for the data of Exercise~1 and test the significance of regression with $\alpha=0.05$. Verify $F_0=t_0^2$ and compute $R^2$.
\end{exercise}

\begin{exercise}
Using Exercise~1, find a $95\%$ confidence interval on the mean response at $x_0=6$ and a $95\%$ prediction interval for a new observation at $x_0=6$. Why is the second wider? What happens at $x_0=20$?
\end{exercise}

\begin{exercise}
Prove that the residuals of a least squares line with an intercept satisfy $\sum e_i=0$ and $\sum x_ie_i=0$, and deduce that the fitted line passes through $(\bar{x},\bar{y})$.
\end{exercise}

\begin{exercise}
Show that $\hat\beta_1=R\,s_y/s_x$ and $R^2=SS_R/SS_T$. If $R=0.8$ and $s_y=10$, $s_x=2$, find $\hat\beta_1$ and the proportion of variance explained.
\end{exercise}

\begin{exercise}
The data below follow an approximate relationship $y=\beta_0e^{\beta_1x}$: $x=1,2,3,4,5$; $y=2.7,7.1,20.5,54.9,148.0$. Linearise, fit by least squares and estimate $\beta_0$ and $\beta_1$.
\end{exercise}

\begin{exercise}
In a sample of $n=27$ pairs $r=-0.45$. Test $H_0:\rho=0$ and find a $95\%$ confidence interval for $\rho$ using Fisher's transformation.
\end{exercise}

\begin{exercise}
Explain, with the help of a residual plot sketch, how to recognise (a) a nonlinear relationship, (b) non-constant variance and (c) a single outlier, and what you would do about each.
\end{exercise}
